% from ega4-9-fr.tex
\oldpage[IV-3]{81}

pour $s\in U$, $w_s:W_s\to X_s$ est une immersion ouverte (9.6.1, (x)), et que
son image dans $X_s$ est partout dense (9.5.3). Autrement dit, pour $s\in U$,
il y a dans $X_s$ un ouvert partout dense $W'_s$ qui est géométriquement
intègre~; le critère (4.5.9, \emph{c)}) montre donc déjà que $X_s$ est
géométriquement irréductible pour $s\in U$. Enfin, comme $X_\eta$ est réduit et
par suite n'a pas de cycle premier associé immergé (3.2.1), on peut aussi
supposer que pour $s\in U$, $X_s$ n'a pas de cycle premier associé immergé
(9.7.6)~; soit alors (pour un $s\in U$ fixé) $\Omega$ une extension
algébriquement close de $\boldsymbol{k}(s)$, et soit
$p:(X_s)_{(\Omega)}\to X_s$ la projection canonique~; $p^{-1}(W'_s)$ est un
ouvert dense dans $(X_s)_{(\Omega)}$ (2.3.10) et est intègre par hypothèse~; en
outre $(X_s)_{(\Omega)}$ n'a pas de cycle premier associé immergé (4.2.7), donc
on conclut de (3.2.1) que $(X_s)_{(\Omega)}$ est réduit~; cela achève de prouver
que $X_s$ est géométriquement intègre (4.6.1).

2$^\circ$ Supposons que $X_\eta$ soit géométriquement irréductible~; comme
$X_{\mathrm{red}}$ est aussi de type fini sur $S$ (1.5.4, (vi)) on peut, en
tenant compte de (\textbf{I}, 5.1.8), remplacer $X$ par $X_{\mathrm{red}}$~;
alors $X_\eta$ est aussi intègre, et comme par hypothèse $X_\eta$ est
géométriquement réduit, il est géométriquement intègre. On est alors dans les
conditions de 1$^\circ$ et on en conclut (revenant aux hypothèses initiales) que
$X_s$ est géométriquement irréductible pour $s$ voisin de $\eta$.

3$^\circ$ Supposons $X_\eta$ géométriquement connexe, et soient $Z_i$
($1\leqslant i\leqslant n$) les composantes irréductibles de $X_\eta$~; il
existe (en vertu de (5.10.8.1) appliqué à $\mathfrak T=\{\varnothing\}$) une
application \emph{surjective} $j\mapsto\nu(j)$ d'un intervalle $[1,m]$ de
$\mathbf N$ sur $[1,n]$ telle que
$Z_{\nu(j)}\cap Z_{\nu(j+1)}\neq\varnothing$ pour $1\leqslant j\leqslant m$.
Pour tout $i$, soit $X_i$ l'adhérence de $Z_i$ dans $X$, et soit $Y$ la réunion
des $X_i$~; comme $Y_\eta=X_\eta$ par définition, on peut supposer, en vertu de
(9.5.1), que $Y_s=X_s$ pour tout $s\in S$, donc que $X_s$ est réunion des
$(X_i)_s$. Or, en considérant les sous-préschémas fermés réduits de $X$ ayant
les $X_i$ pour espaces sous-jacents, on voit par 2$^\circ$ qu'il existe un
voisinage $U$ de $\eta$ dans $S$ tel que pour $s\in U$ les $(X_i)_s$ soient
\emph{géométriquement irréductibles} (puisque les $(X_i)_\eta$ peuvent être
supposés géométriquement irréductibles, comme on l'a vu au début). En outre, on
peut aussi supposer que pour $s\in U$, on a
$(X_{\nu(j)})_s\cap(X_{\nu(j+1)})_s\neq\varnothing$ (9.5.1) pour
$1\leqslant j\leqslant m$~; on en conclut aussitôt que $X_s$ est connexe, donc
(4.5.13.1) géométriquement connexe pour $s\in U$.

4$^\circ$ Supposons $X_\eta$ géométriquement réduit~; soient $Z_i$ les
composantes irréductibles de $X_\eta$, $W_i$ l'intérieur de $Z_i$ dans
$X_\eta$~; il y a pour chaque $i$ un ouvert $V_i$ de $X$ tel que
$W_i=V_i\cap X_\eta$ pour tout $i$~; comme les $W_i$ sont ouverts et deux à deux
disjoints et que leur réunion est dense dans $X_\eta$, on peut (9.5.1, 9.5.3 et
9.5.4) supposer que pour $s$ voisin de $\eta$, les $(V_i)_s$ sont des ouverts
deux à deux disjoints dans $X_s$ et que leur réunion est dense dans $X_s$. En
outre, comme les $W_i$ sont géométriquement réduits et ont été supposés au
début géométriquement irréductibles, il résulte du 1$^\circ$ que pour $s$ voisin
de $\eta$, les $(V_i)_s$ sont géométriquement intègres, et \emph{a fortiori}
réduits. D'autre part, on tire de (9.7.6) que pour $s$ voisin de $\eta$, $X_s$
n'a pas de cycle premier associé immergé, puisqu'il en est ainsi de $X_\eta$
qui est réduit (3.2.1)~; on conclut donc de (3.2.1) que $X_s$ est réduit, et de
(4.6.1) qu'il est géométriquement réduit.

